\BeleskeID{03-s021}
% element: 03-s021:e01
% element: 03-s021:e02
% element: 03-s021:e03
% element: 03-s021:e04
% element: 03-s021:d01

Najviša grupa dobija \(EI=1\). Svaka sledeća preuzima EI sa EO prethodne grupe, pa dozvola prolazi samo kroz grupe bez zahteva. Grupa sa najvišim aktivnim ulazom zadržava dozvolu, koduje lokalni indeks i svojim EO zaustavlja sve niže grupe. Njihovi adresni bitovi tada su nula, pa dve četvoroulazne ILI grane propuštaju niže bitove samo pobedničke grupe.

Na primer, za aktivne \(I_{10},I_5,I_1\), grupa 3 nema zahtev i propušta dozvolu grupi 2. U grupi 2 pobednik je lokalni ulaz 2, pa su niži bitovi 10. Dozvola za grupe 1 i 0 je nula; drugi nivo iz GS signala koduje grupu 2 kao 10. Ukupni kod je 1010, odnosno indeks 10. Ako nema zahteva, svi adresni bitovi i GS su nula, a dozvola prolazi kroz ceo lanac.

Visoka impedansa znači da izlaz ne nameće ni nulu ni jedinicu zajedničkom vodu. Direktno objedinjavanje bilo bi moguće sa izlazima koji zaista podržavaju to stanje, uz uslov da u datom trenutku samo izabrana grupa aktivno upravlja vodom, a sve ostale su u visokoj impedansi. Ako su svi izlazi u tom stanju, vod sam po sebi nema garantovan nivo; pri \(GS=0\) njegovu vrednost ne treba tumačiti kao validnu adresu. Običan neaktivan izlaz na logičkoj nuli nije izlaz visoke impedanse.

Za simbole kola srednjeg stepena integracije ne koristi se jedan strogo jedinstven raspored priključaka. Uobičajeno je razdvajanje ulaza i izlaza, često levo i desno; bitno je da nijedan priključak ne ostane bez oznake. Tako se može pratiti koji vod nosi podatak, adresu, dozvolu rada ili informaciju o prisustvu zahteva.
